\documentclass[10pt,a4paper]{amsart}
\usepackage{amsmath,amssymb,amsthm,graphicx}
\usepackage{verbatim}
\newenvironment{reference}{\comment}{\endcomment}
%\newenvironment{reference}{}{}
\newenvironment{slogan}{\comment}{\endcomment}
\newenvironment{history}{\comment}{\endcomment}

% For commutative diagrams we use Xy-pic
\usepackage[all]{xy}

% We use 2cell for 2-commutative diagrams.
\xyoption{2cell}
% Theorem environments.
%
\theoremstyle{plain}
\newtheorem{theorem}[subsection]{Theorem}
\newtheorem{proposition}[subsection]{Proposition}
\newtheorem{lemma}[subsection]{Lemma}

\theoremstyle{definition}
\newtheorem{definition}[subsection]{Definition}
\newtheorem{example}[subsection]{Example}
\newtheorem{exercise}[subsection]{Exercise}
\newtheorem{situation}[subsection]{Situation}

\theoremstyle{remark}
\newtheorem{remark}[subsection]{Remark}
\newtheorem{remarks}[subsection]{Remarks}

\numberwithin{equation}{subsection}

% Macros
%
\def\lim{\mathop{\mathrm{lim}}\nolimits}
\def\colim{\mathop{\mathrm{colim}}\nolimits}
\def\Spec{\mathop{\mathrm{Spec}}}
\def\Hom{\mathop{\mathrm{Hom}}\nolimits}
\def\Ext{\mathop{\mathrm{Ext}}\nolimits}
\def\SheafHom{\mathop{\mathcal{H}\!\mathit{om}}\nolimits}
\def\SheafExt{\mathop{\mathcal{E}\!\mathit{xt}}\nolimits}
\def\Sch{\mathit{Sch}}
\def\Mor{\mathop{\mathrm{Mor}}\nolimits}
\def\Ob{\mathop{\mathrm{Ob}}\nolimits}
\def\Sh{\mathop{\mathit{Sh}}\nolimits}
\def\NL{\mathop{N\!L}\nolimits}
\def\CH{\mathop{\mathrm{CH}}\nolimits}
\def\proetale{{pro\text{-}\acute{e}tale}}
\def\etale{{\acute{e}tale}}
\def\QCoh{\mathit{QCoh}}
\def\Ker{\mathop{\mathrm{Ker}}}
\def\Im{\mathop{\mathrm{Im}}}
\def\Coker{\mathop{\mathrm{Coker}}}
\def\Coim{\mathop{\mathrm{Coim}}}

% Boxtimes
%
\DeclareMathSymbol{\boxtimes}{\mathbin}{AMSa}{"02}

%
% Macros for moduli stacks/spaces
%
\def\QCohstack{\mathcal{QC}\!\mathit{oh}}
\def\Cohstack{\mathcal{C}\!\mathit{oh}}
\def\Spacesstack{\mathcal{S}\!\mathit{paces}}
\def\Quotfunctor{\mathrm{Quot}}
\def\Hilbfunctor{\mathrm{Hilb}}
\def\Curvesstack{\mathcal{C}\!\mathit{urves}}
\def\Polarizedstack{\mathcal{P}\!\mathit{olarized}}
\def\Complexesstack{\mathcal{C}\!\mathit{omplexes}}
% \Pic is the operator that assigns to X its picard group, usage \Pic(X)
% \Picardstack_{X/B} denotes the Picard stack of X over B
% \Picardfunctor_{X/B} denotes the Picard functor of X over B
\def\Pic{\mathop{\mathrm{Pic}}\nolimits}
\def\Picardstack{\mathcal{P}\!\mathit{ic}}
\def\Picardfunctor{\mathrm{Pic}}
\def\Deformationcategory{\mathcal{D}\!\mathit{ef}}
\usepackage{cleveref}
\crefname{originaldefinition}{original definition}{original definitions}
\Crefname{originaldefinition}{Original definition}{Original definitions}
\makeatletter
\newlabel{decent-spaces-definition-relative-conditions}{{03KZ}{source}{}{sourcepointer-1}{}}
\newlabel{decent-spaces-definition-relative-conditions@cref}{{[originaldefinition][0][]03KZ}{[1][0][]source}{}{}{}}
\makeatother
\crefname{originaldefinition}{original definition}{original definitions}
\Crefname{originaldefinition}{Original definition}{Original definitions}
\makeatletter
\newlabel{decent-spaces-definition-very-reasonable}{{03I8}{source}{}{sourcepointer-2}{}}
\newlabel{decent-spaces-definition-very-reasonable@cref}{{[originaldefinition][0][]03I8}{[1][0][]source}{}{}{}}
\makeatother
\crefname{originallemma}{original lemma}{original lemmas}
\Crefname{originallemma}{Original lemma}{Original lemmas}
\makeatletter
\newlabel{decent-spaces-lemma-decent-specialization}{{03IL}{source}{}{sourcepointer-3}{}}
\newlabel{decent-spaces-lemma-decent-specialization@cref}{{[originallemma][0][]03IL}{[1][0][]source}{}{}{}}
\makeatother
\crefname{originallemma}{original lemma}{original lemmas}
\Crefname{originallemma}{Original lemma}{Original lemmas}
\makeatletter
\newlabel{decent-spaces-lemma-property-for-morphism-out-of-property}{{0ABX}{source}{}{sourcepointer-4}{}}
\newlabel{decent-spaces-lemma-property-for-morphism-out-of-property@cref}{{[originallemma][0][]0ABX}{[1][0][]source}{}{}{}}
\makeatother
\crefname{originallemma}{original lemma}{original lemmas}
\Crefname{originallemma}{Original lemma}{Original lemmas}
\makeatletter
\newlabel{decent-spaces-lemma-topology-fibre}{{0AC8}{source}{}{sourcepointer-5}{}}
\newlabel{decent-spaces-lemma-topology-fibre@cref}{{[originallemma][0][]0AC8}{[1][0][]source}{}{}{}}
\makeatother
\crefname{originaltheorem}{original theorem}{original theorems}
\Crefname{originaltheorem}{Original theorem}{Original theorems}
\makeatletter
\newlabel{flat-theorem-check-flatness-at-associated-points}{{05UA}{source}{}{sourcepointer-6}{}}
\newlabel{flat-theorem-check-flatness-at-associated-points@cref}{{[originaltheorem][0][]05UA}{[1][0][]source}{}{}{}}
\makeatother
\title{Flatness at a point can be checked at the associated points of its fibre}
\author{The Stacks Project Authors}
\date{stacks-project:spaces-flat.tex:local-snapshot-20261012; GNU FDL 1.2 or later}
\begin{document}
\maketitle
\noindent\textit{Editorial note.} All delivered mathematical text is the author's own, copied from the pinned Stacks Project chapter file spaces-flat.tex; no statement or proof step was rewritten, reordered or removed.\par\medskip
\noindent\textit{Editorial note.} The counted proof is the author's complete argument: an etale chart reduction with a comparison of associated points and of specializations, and a decentness argument for the converse direction\par\medskip
\noindent\textit{Editorial note.} The label theorem-check-flatness-at-associated-points names the counted statement, and the proof environment immediately following it in the source is delivered verbatim as its proof.\par\medskip
\noindent\textit{Editorial note.} The counted statement is carried by the author's own theorem environment in the chapter source.\par\medskip
\noindent\textit{Allowed standard prerequisites (editorial).}\begin{itemize}
\item The associated points of the fibres of a finite type quasicoherent module, and their behaviour under etale base change
\item Etale charts of a morphism and continuity of the induced map on underlying topological spaces, so that specializations can be compared
\item The notion of a decent algebraic space and of a decent morphism, the representation of a point of a decent space by a quasi-compact monomorphism, and the lifting of specializations along a decent morphism
\item The corresponding criterion in the case of schemes
\end{itemize}\medskip
\ifdefined\hypertarget\hypertarget{sourcepointer-1}{}\fi
\noindent\textit{Source reference.} Original definition 03KZ: sources/stacks-chow/tags.txt, source lines 4695--4695. Prerequisite pointer invoked at proof line(s) 878; it names a standard fact of a sibling chapter that the counted argument uses, and every step of the counted proof built on it is retained verbatim in the delivered body.\par\medskip
\ifdefined\hypertarget\hypertarget{sourcepointer-2}{}\fi
\noindent\textit{Source reference.} Original definition 03I8: sources/stacks-chow/tags.txt, source lines 4594--4594. Prerequisite pointer invoked at proof line(s) 874; it names a standard fact of a sibling chapter that the counted argument uses, and every step of the counted proof built on it is retained verbatim in the delivered body.\par\medskip
\ifdefined\hypertarget\hypertarget{sourcepointer-3}{}\fi
\noindent\textit{Source reference.} Original lemma 03IL: sources/stacks-chow/tags.txt, source lines 4607--4607. Prerequisite pointer invoked at proof line(s) 884; it names a standard fact of a sibling chapter that the counted argument uses, and every step of the counted proof built on it is retained verbatim in the delivered body.\par\medskip
\ifdefined\hypertarget\hypertarget{sourcepointer-4}{}\fi
\noindent\textit{Source reference.} Original lemma 0ABX: sources/stacks-chow/tags.txt, source lines 13035--13035. Prerequisite pointer invoked at proof line(s) 880; it names a standard fact of a sibling chapter that the counted argument uses, and every step of the counted proof built on it is retained verbatim in the delivered body.\par\medskip
\ifdefined\hypertarget\hypertarget{sourcepointer-5}{}\fi
\noindent\textit{Source reference.} Original lemma 0AC8: sources/stacks-chow/tags.txt, source lines 13046--13046. Prerequisite pointer invoked at proof line(s) 876; it names a standard fact of a sibling chapter that the counted argument uses, and every step of the counted proof built on it is retained verbatim in the delivered body.\par\medskip
\ifdefined\hypertarget\hypertarget{sourcepointer-6}{}\fi
\noindent\textit{Source reference.} Original theorem 05UA: sources/stacks-chow/tags.txt, source lines 7509--7509. Prerequisite pointer invoked at proof line(s) 867; it names a standard fact of a sibling chapter that the counted argument uses, and every step of the counted proof built on it is retained verbatim in the delivered body.\par\medskip
\begin{theorem}
\label{theorem-check-flatness-at-associated-points}
Let $S$ be a scheme.
Let $f : X \to Y$ be a morphism of algebraic spaces over $S$
which is locally of finite type.
Let $\mathcal{F}$ be a quasi-coherent $\mathcal{O}_X$-module of finite type.
Let $x \in |X|$ with image $y \in |Y|$.
Set $F = f^{-1}(\{y\}) \subset |X|$.
Consider the conditions
\begin{enumerate}
\item $\mathcal{F}$ is flat at $x$ over $Y$, and
\item for every $x' \in F \cap \text{Ass}_{X/Y}(\mathcal{F})$ which
specializes to $x$ we have that $\mathcal{F}$ is flat at $x'$ over $Y$.
\end{enumerate}
Then we always have (2) $\Rightarrow$ (1). If $X$ and $Y$ are
decent, then (1) $\Rightarrow$ (2).
\end{theorem}

\begin{proof}
Assume (2).
Choose a scheme $V$ and a surjective \'etale morphism $V \to Y$.
Choose a scheme $U$ and a surjective \'etale morphism $U \to V \times_Y X$.
Choose a point $u \in U$ mapping to $x$. Let $v \in V$ be the image of $u$.
We will deduce the result from the corresponding result for
$\mathcal{F}|_U = (U \to X)^*\mathcal{F}$ and the point $u$.
$U_v$. This works because $\text{Ass}_{U/V}(\mathcal{F}|_U) \cap |U_v|$
is equal to $\text{Ass}_{U_v}(\mathcal{F}|_{U_v})$ and equal to the inverse
image of $F \cap \text{Ass}_{X/Y}(\mathcal{F})$.
Since the map $|U_v| \to F$ is continuous we see that
specializations in $|U_v|$ map to specializations in $F$,
hence condition (2) is inherited by $U \to V$,
$\mathcal{F}|_U$, and the point $u$.
Thus More on Flatness, Theorem
\ref{flat-theorem-check-flatness-at-associated-points} applies
and we conclude that (1) holds.

\medskip\noindent
If $Y$ is decent, then we can represent
$y$ by a quasi-compact monomorphism $\Spec(k) \to Y$
(by definition of decent spaces, see
Decent Spaces, Definition \ref{decent-spaces-definition-very-reasonable}).
Then $F = |X_k|$, see
Decent Spaces, Lemma \ref{decent-spaces-lemma-topology-fibre}.
If in addition $X$ is decent (or more generally if $f$ is decent, see
Decent Spaces, Definition \ref{decent-spaces-definition-relative-conditions}
and Decent Spaces, Lemma
\ref{decent-spaces-lemma-property-for-morphism-out-of-property}),
then $X_y$ is a decent space too. Furthermore, specializations in
$F$ can be lifted to specializations
in $U_v \to X_y$, see
Decent Spaces, Lemma \ref{decent-spaces-lemma-decent-specialization}.
Having said this it is clear that the reverse implication
holds, because it holds in the case of schemes.
\end{proof}

\end{document}
